\documentclass[11pt,oneside,openany]{book}
\usepackage[T1]{fontenc}
\usepackage[utf8]{inputenc}
\usepackage[sc]{mathpazo}
\renewcommand{\sfdefault}{lmss}
\renewcommand{\ttdefault}{lmtt}
\usepackage{amsmath,amssymb,amsthm,mathtools}
\usepackage[paperwidth=7.5in,paperheight=10in,inner=.85in,outer=.85in,top=.85in,bottom=.85in,headheight=15pt,headsep=20pt]{geometry}
\usepackage{microtype}
\usepackage{xcolor,graphicx,booktabs,array,longtable}
\newcolumntype{L}[1]{>{\raggedright\arraybackslash}p{#1}}
\usepackage{tikz,pgfplots}
\usetikzlibrary{arrows.meta,calc,patterns,decorations.markings}
\pgfplotsset{compat=1.18}
\definecolor{ink}{HTML}{17334B}
\definecolor{teal}{HTML}{137C83}
\definecolor{copper}{HTML}{AE5932}
\definecolor{soft}{HTML}{EEF3F5}
\pgfplotsset{every axis/.append style={axis line style={black!65},tick style={black!60},tick label style={font=\small},label style={font=\small},legend style={font=\small,draw=none,fill=white},grid style={black!10},major grid style={black!10},line width=.6pt},every axis plot/.append style={line width=1.15pt}}
\usepackage{enumitem,needspace,etoolbox}
\setlist[enumerate]{topsep=5pt,itemsep=5pt,parsep=0pt}
\setlist[itemize]{topsep=4pt,itemsep=3pt,parsep=0pt}
\usepackage{titlesec,fancyhdr}
\titleformat{\chapter}[display]{\normalfont\color{ink}}{\large\sffamily\MakeUppercase{Chapter \thechapter}}{12pt}{\huge\bfseries\raggedright\hyphenpenalty=10000\exhyphenpenalty=10000}[\vspace{12pt}\titlerule]
\titlespacing*{\chapter}{0pt}{-10pt}{24pt}
\titleformat{\section}{\large\bfseries\color{ink}}{\thesection}{.65em}{}
\titleformat{\subsection}{\normalsize\bfseries\color{ink}}{\thesubsection}{.6em}{}
\titlespacing*{\section}{0pt}{18pt plus 3pt minus 2pt}{7pt}
\titlespacing*{\subsection}{0pt}{12pt plus 2pt minus 2pt}{5pt}
\pagestyle{fancy}
\fancyhf{}
\fancyhead[L]{\small\nouppercase{\leftmark}}
\fancyfoot[C]{\small\thepage}
\renewcommand{\headrulewidth}{.3pt}
\renewcommand{\chaptermark}[1]{\markboth{\thechapter\quad #1}{}}
\fancypagestyle{plain}{\fancyhf{}\fancyfoot[C]{\small\thepage}\renewcommand{\headrulewidth}{0pt}}
\usepackage[font=small,labelfont=bf,labelsep=period]{caption}
\usepackage{float,placeins}
\floatplacement{figure}{htbp}
\setcounter{topnumber}{3}
\setcounter{bottomnumber}{2}
\renewcommand{\topfraction}{.9}
\renewcommand{\textfraction}{.08}
\renewcommand{\floatpagefraction}{.8}
\usepackage{makeidx}
\makeindex
\usepackage[pdfencoding=auto,psdextra,colorlinks=true,linkcolor=ink,citecolor=teal,urlcolor=teal,bookmarksnumbered=true]{hyperref}
\usepackage{bookmark}
\hypersetup{pdftitle={Discrete Information Geometry},pdfsubject={Finite differences, probability, and information flows}}
\urlstyle{same}
\newtheorem{theorem}{Theorem}[chapter]
\newtheorem{proposition}[theorem]{Proposition}
\newtheorem{lemma}[theorem]{Lemma}
\newtheorem{corollary}[theorem]{Corollary}
\theoremstyle{definition}
\newtheorem{definition}[theorem]{Definition}
\newtheorem{example}{Example}[chapter]
\BeforeBeginEnvironment{example}{\par\addvspace{4pt}\noindent\begin{minipage}{\linewidth}}
\AfterEndEnvironment{example}{\end{minipage}\par\addvspace{4pt}}
\theoremstyle{remark}
\newtheorem*{remark}{Remark}
\newenvironment{solution}{\par\smallskip\noindent\textit{Solution.}\ }{\par\medskip}
\newenvironment{exercises}{\begin{enumerate}[label=\thechapter.\arabic*,ref=\thechapter.\arabic*,leftmargin=2.8em,itemsep=6pt]}{\end{enumerate}}
\newenvironment{answers}{\begin{description}[leftmargin=2.9em,labelwidth=2.5em,labelsep=.4em,itemsep=8pt,font=\normalfont\bfseries]}{\end{description}}
\newcommand{\R}{\mathbb{R}}
\newcommand{\E}{\mathbb{E}}
\newcommand{\Var}{\operatorname{Var}}
\newcommand{\Cov}{\operatorname{Cov}}
\newcommand{\KL}[2]{D\!\left(#1\,\middle\|\,#2\right)}
\newcommand{\grad}{\operatorname{grad}}
\newcommand{\arcosh}{\operatorname{arcosh}}
\newcommand{\sech}{\operatorname{sech}}
\newcommand{\softmax}{\operatorname{softmax}}
\newcommand{\diag}{\operatorname{diag}}
\newcommand{\dd}{\,\mathrm{d}}
\newcommand{\ip}[2]{\left\langle #1,#2\right\rangle}
\newcommand{\p}{\mathbf{p}}
\newcommand{\q}{\mathbf{q}}
\newcommand{\f}{\mathbf{f}}
\newcommand{\uvec}{\mathbf{u}}
\newcommand{\vvec}{\mathbf{v}}
\newcommand{\ones}{\mathbf{1}}
\newcommand{\simp}{\Delta^{n-1}}
\newcommand{\simpi}{\Delta^{n-1}_{\!\circ}}
\numberwithin{equation}{chapter}
\setlength{\parindent}{1.2em}
\setlength{\parskip}{0pt}
\linespread{1.07}
\setlength{\emergencystretch}{2em}
\clubpenalty=10000
\widowpenalty=10000
\displaywidowpenalty=10000
\raggedbottom
\setcounter{tocdepth}{1}
\makeatletter
\renewcommand*\l@section{\@dottedtocline{1}{1.5em}{3.2em}}
\makeatother
\setcounter{secnumdepth}{2}
\providecommand{\tightlist}{\setlength{\itemsep}{0pt}\setlength{\parskip}{0pt}}
\providecommand{\passthrough}[1]{#1}
\newcommand{\labheading}{\FloatBarrier\section*{Laboratory investigation}\addcontentsline{toc}{section}{Laboratory investigation}}

\begin{document}
\begin{titlepage}\color{ink}\centering\vspace*{1.0in}{\small\sffamily PROBABILITY AND MATHEMATICAL MODELLING\par}\vspace{.45in}{\Huge\bfseries Discrete Information\\[.12in]Geometry\par}\vspace{.35in}{\Large Finite Differences, Probability,\\and Information Flows\par}\vspace{1cm}{\Large Chapter 8\par}\vfill\begin{tikzpicture}[scale=1.05]\draw[teal,line width=1pt](0,0)--(4,0)--(2,3.464)--cycle;\foreach \a in {0,...,6}{\foreach \b in {0,...,6}{\pgfmathtruncatemacro{\c}{6-\a-\b}\ifnum\c>-1\fill[teal] ({(\b+\c/2)*4/6},{\c*3.464/6}) circle(1.5pt);\fi}}\end{tikzpicture}\vfill{\large A course in ten chapters\par}\vspace{.3in}{\small\url{https://bobovski66.github.io/IntroDiscreteInfoGeom/}\par}\end{titlepage}
\setcounter{chapter}{7}

\chapter{Path sums, circulation, and statistical geodesics}
A distribution can move from one state to another in many ways. A survey composition can be changed one response at a time. A model can be adjusted by mixing it with a target. A sequence of distributions can also follow equal increments of statistical distance. The endpoints alone do not specify the permitted journey.

Chapter 7 gave us an exact distance between distributions. We now distinguish that endpoint distance from quantities accumulated along a route. The discrete fundamental theorem will identify sums that depend only on endpoints. Positive edge lengths will lead to shortest-path problems, and the sphere construction will give an explicit shortest route when intermediate states are unrestricted.

\section{A path is a finite sequence of legal moves}\label{ch08-a-path-is-a-finite-sequence-of-legal-moves}

Let a graph have probability states as its vertices. Two states are connected when the model permits a direct move between them. A \textbf{path} is a sequence \[\gamma=(v_0,v_1,\ldots,v_m)\] in which every consecutive pair is connected. The number \(m\) counts moves, while \(m+1\) counts listed states.

Assign a signed number \(A(u,v)\) to a directed edge. Require \[A(v,u)=-A(u,v).\] This \textbf{antisymmetry} says reversing a step reverses its contribution. The discrete line integral is the finite sum \[\boxed{\int_\gamma^d A=\sum_{j=0}^{m-1}A(v_j,v_{j+1}).}\] The integral symbol is a name for this operation on edges. Its definition is ordinary finite addition.

Such an edge value can represent a change of potential, a signed work contribution, or another quantity whose sign reverses with direction. A general travel cost need not have this property and will be handled separately.

\section{Gradient integrals recover an endpoint change}\label{ch08-gradient-integrals-recover-an-endpoint-change}

If a scalar potential \(\phi\) is known at vertices, define \[A_\phi(u,v)=\phi(v)-\phi(u).\] Along any path, \[\begin{aligned}
\int_\gamma^d A_\phi
&=[\phi(v_1)-\phi(v_0)]+\cdots+[\phi(v_m)-\phi(v_{m-1})]\\
&=\phi(v_m)-\phi(v_0).
\end{aligned}\] This is the \textbf{discrete fundamental theorem for line integrals}. It is the telescoping identity from Chapter 1 applied to values encountered along a path.

\begin{example}[Different routes, identical potential change]

Suppose \(\phi(U)=0\), \(\phi(V)=2\), \(\phi(W)=5\), and \(\phi(Z)=3\). Along \(U,V,W\), the contributions are \(2\) and \(3\), totaling five. Along \(U,Z,W\), they are \(3\) and \(2\), again totaling five.

The route \(U,W,V,W\) contributes \(5,-3,3\), still totaling five. Backtracking adds movement, but its signed potential changes cancel.

\end{example}

If \(\phi\) is entropy, the integral gives the net entropy change. It does not report the largest entropy encountered or the total size of all fluctuations. Choosing the quantity being accumulated is essential.

Around a closed path, the endpoint is the starting point, so every gradient integral is zero. Chapter 3 used this fact as a test for a gradient field.

\section{Circulation tests whether a potential exists}\label{ch08-circulation-tests-whether-a-potential-exists}

Consider a triangle oriented \(U\to V\to W\to U\), with edge values all equal to one in that direction. Its circulation is three. If these values came from a potential, adding them would telescope to zero. Thus no potential can generate this field.

There is a converse on a connected graph: an antisymmetric edge field is a gradient exactly when every closed-path sum is zero.

To construct the potential, choose a root vertex \(r\) and assign \(\phi(r)=0\). For any other vertex \(v\), choose a path from \(r\) to \(v\) and define \(\phi(v)\) to be its edge sum. This definition must be independent of the chosen path. Two candidate paths followed in opposite directions form a closed route. Its sum is zero by assumption, so the two sums agree.

Finally, follow a path to \(u\) and then the edge \(u\to v\). The construction gives \[\phi(v)=\phi(u)+A(u,v),\] hence \(A(u,v)=\phi(v)-\phi(u)\). This proves the claim.

For a finite graph, one can choose a spanning tree: a connected collection of edges containing every vertex and no cycles. Every remaining edge creates one basic cycle with the tree. Checking those cycle sums suffices, because the tree paths already define a candidate potential. This is a finite way to test the global condition.

There is a useful local-to-boundary version of the circulation calculation. Suppose a patch of triangular lattice faces is assigned a consistent counterclockwise orientation. Add the circulation around every face. An edge shared by two faces is traversed in opposite directions, so its two contributions cancel. Only the outer boundary edges remain: \[\boxed{\text{sum of oriented face circulations}
=\text{circulation around the patch boundary}.}\] This is a discrete Stokes identity. It is another exact cancellation theorem, now relating local loops to a larger loop.

For two triangles sharing an edge, one can check it by writing out both three-edge sums. The shared edge appears once as \(A(u,v)\) and once as \(A(v,u)=-A(u,v)\). The remaining four terms form the boundary route. A larger patch works by the same edge-by-edge accounting.

On a triangle patch without holes, zero circulation on every small face forces zero circulation on its boundary loops. On a graph with a hole, a loop around the hole may not bound a collection of included faces. Local face checks then need not detect its circulation. The earlier criterion used every closed path precisely because it does not assume the graph has a particular filling.

\section{Path length accumulates positive costs}\label{ch08-path-length-accumulates-positive-costs}

Give each unoriented edge a positive weight \(\ell(u,v)=\ell(v,u)\). Define \[L(\gamma)=\sum_{j=0}^{m-1}\ell(v_j,v_{j+1}).\] These contributions stay positive when a step is reversed. A round trip therefore has positive length even when its gradient integral is zero.

\begin{longtable}[]{@{}
  >{\raggedright\arraybackslash}p{(\columnwidth - 4\tabcolsep) * \real{0.3333}}
  >{\raggedright\arraybackslash}p{(\columnwidth - 4\tabcolsep) * \real{0.3333}}
  >{\raggedright\arraybackslash}p{(\columnwidth - 4\tabcolsep) * \real{0.3333}}@{}}
\toprule\noalign{}
\begin{minipage}[b]{\linewidth}\raggedright
Quantity along a path
\end{minipage} & \begin{minipage}[b]{\linewidth}\raggedright
Edge contribution
\end{minipage} & \begin{minipage}[b]{\linewidth}\raggedright
Effect of reversing a step
\end{minipage} \\
\midrule\noalign{}
\endhead
\bottomrule\noalign{}
\endlastfoot
Potential change & \(\phi(v)-\phi(u)\) & Sign reverses \\
Statistical length & \(d_{\mathrm{FR}}(u,v)\) & Length is unchanged \\
Transfer count & 1 per allowed transfer & Another move is counted \\
\end{longtable}

A shortest path minimizes the sum of the chosen costs among permitted paths with the same endpoints. Before solving such a problem, specify both the allowed moves and their weights.

For statistical edge weights, the triangle inequality gives \[L(\gamma)\geq d_{\mathrm{FR}}(v_0,v_m).\] Repeatedly apply the triangle inequality to insert the intermediate states. The endpoint distance is a lower bound on every route's statistical length.

\section{Token distance can be solved by counting surpluses}\label{ch08-token-distance-can-be-solved-by-counting-surpluses}

Let count vectors \(a\) and \(b\) each have total \(N\), and permit a transfer between any two categories. Each transfer decreases one count by one and increases another by one.

The total surplus relative to the target is \[S=\sum_{i:a_i>b_i}(a_i-b_i).\] Since the total counts agree, the total deficit is also \(S\). Thus \[\boxed{d_{\mathrm{token}}(a,b)=S=\frac12\sum_i|a_i-b_i|.}\] At least \(S\) transfers are needed because each transfer can correct at most one surplus token. A route attaining the bound repeatedly sends a token from any surplus category to any deficit category. The donor has a token available, so every move is legal.

For \((4,1,1)\to(1,2,3)\), the differences are \((3,-1,-2)\). Three transfers suffice: one from \(A\) to \(B\) and two from \(A\) to \(C\), in any order. Their statistical lengths depend on the intermediate compositions even though their transfer counts agree.

The formula assumes every category-to-category transfer is allowed. A network restriction can force additional intermediate moves and change the answer.

\section{A lattice restriction can produce a measurable detour}\label{ch08-a-lattice-restriction-can-produce-a-measurable-detour}

Consider the two-token, three-category lattice. Start at \(p=(1,0,0)\) and end at \(r=(0,1/2,1/2)\). Their supports are disjoint, so the direct spherical distance is \(\pi\).

Two legal transfers produce the route \[p\ \longrightarrow\ q=(1/2,1/2,0)\ \longrightarrow\ r.\] The first overlap is \(1/\sqrt2\), giving distance \(\pi/2\). The second overlap is \(1/2\), giving distance \(2\pi/3\). The route length is therefore \[\frac\pi2+\frac{2\pi}3=\frac{7\pi}6.\] Starting instead through \((1/2,0,1/2)\) gives the same length.

These are the only possible first neighbors of \(p\). After either first step, any remaining route has length at least \(2\pi/3\) by the triangle inequality. Hence \(7\pi/6\) is the shortest lattice-path length. Its excess over the unrestricted endpoint distance is \(\pi/6\).

This gap measures a concrete constraint. The unrestricted shortest arc passes through distributions unavailable to a two-token allocation.

\section{Finding weighted shortest paths systematically}\label{ch08-finding-weighted-shortest-paths-systematically}

For a finite graph with nonnegative edge weights, \textbf{Dijkstra's algorithm} computes shortest distances from a starting vertex.

Give the start tentative distance zero and every other vertex tentative distance infinity. At each round, select an unvisited vertex with the smallest tentative distance. For each neighbor, compare its current tentative distance with the selected vertex's distance plus the connecting edge weight. Replace the neighbor's value if this route is shorter, and record its predecessor. Mark the selected vertex visited.

A visited distance can be finalized because an alternative route that reaches it through an unvisited vertex must already have cost at least the smallest remaining tentative distance. Adding nonnegative edge costs cannot reduce that cost. Stop when the destination is finalized, then follow predecessor records backward to recover a shortest route.

\begin{example}[Following the tentative distances]

Suppose edges \(AB,AC,CB,BD,CD\) have weights \(4,1,2,1,6\). Starting at \(A\):

\begin{longtable}[]{@{}llll@{}}
\toprule\noalign{}
Vertex just finalized & Tentative distance to \(B\) & To \(C\) & To \(D\) \\
\midrule\noalign{}
\endhead
\bottomrule\noalign{}
\endlastfoot
\(A\) & 4 & 1 & \(\infty\) \\
\(C\) & 3 & 1 & 7 \\
\(B\) & 3 & 1 & 4 \\
\(D\) & 3 & 1 & 4 \\
\end{longtable}

The predecessor chain gives \(A,C,B,D\), of total length four. Taking the apparently direct edge \(AB\) would have cost four before even reaching \(D\).

\end{example}

On a token lattice, the vertices are count states and the weights can be exact Fisher--Rao distances. The same algorithm applies. A shortest weighted route answers a finite optimization question; it need not use the fewest transfers if a longer sequence of transfers has lower total weight.

The shortest-path algorithm distinguishes a discovered route from a finalized one. A tentative value is an upper bound because an actual path realizing it has been found. It may fall when a better route is discovered. Once the smallest tentative vertex is finalized, nonnegative weights establish that its upper bound is also a lower bound.

Ties do not invalidate the distance calculation. They can be resolved in any fixed order. A predecessor record may then choose one of several equally short paths. If the purpose is to list all shortest paths, more than one predecessor may need to be retained. The demonstrations need only produce one valid minimizing route.

A zero-weight edge can be handled by the same algorithm. Negative weights undermine the finalization argument because a later edge could reduce an accumulated cost. Fisher--Rao weights between distinct distributions are positive, so the probability lattice falls within the intended setting.

For the two-token detour example, the minimizing route has only one intermediate lattice state. Its unrestricted spherical midpoint is \((1/2,1/4,1/4)\), unavailable with two tokens. With four tokens this midpoint is available, but the direct connections to it are still constrained to one-token transfers. Admitting a desired intermediate state and admitting the needed moves are separate requirements.

\section{Two ways to interpolate between distributions}\label{ch08-two-ways-to-interpolate-between-distributions}

An \textbf{interpolation} specifies intermediate states between endpoints. The simplest is the ordinary mixture \[m(t)=(1-t)p+tr,\qquad 0\leq t\leq1.\] Every coordinate is nonnegative and the total remains one. At \(t=1/2\), it gives the coordinate average. Sampling \(t\) at \(0,1/M,\ldots,1\) produces equal coordinate increments.

To obtain equal statistical increments, work on the sphere. Let \[s=2\sqrt p,\qquad u=2\sqrt r,\qquad
\alpha=\arccos B(p,r).\] The square root is taken coordinate by coordinate. For different endpoints, define \[s(t)=\frac{\sin((1-t)\alpha)}{\sin\alpha}s
+\frac{\sin(t\alpha)}{\sin\alpha}u,\qquad
g_i(t)=\frac{s_i(t)^2}{4}.\] For identical endpoints, set \(g(t)=p\) to avoid dividing by zero. This is \textbf{spherical interpolation}.

The coefficients are nonnegative because their angles lie between zero and \(\pi/2\). The formula therefore stays in the nonnegative part of the sphere. To verify its length, choose orthogonal unit vectors \[e_1=s/2,\qquad
e_2=\frac{u/2-(\cos\alpha)e_1}{\sin\alpha}.\] The dot-product formula shows \(e_1\cdot e_2=0\) and \(\|e_2\|=1\). The sine subtraction identity then rewrites the interpolation as \[s(t)=2[\cos(t\alpha)e_1+\sin(t\alpha)e_2].\] Its squared length is four. Thus the recovered probabilities really add to one.

This form also shows that the angle between the states at \(t_1\) and \(t_2\) is \(|t_2-t_1|\alpha\). Consequently \[d_{\mathrm{FR}}(g(t_1),g(t_2))=2|t_2-t_1|\alpha.\]

\section{Equal statistical steps follow a geodesic}\label{ch08-equal-statistical-steps-follow-a-geodesic}

Choose an integer \(M>0\) and list \(g(0),g(1/M),\ldots,g(1)\). Consecutive distances are \(2\alpha/M\), so their sum is \(2\alpha=d_{\mathrm{FR}}(p,r)\). This route attains the endpoint lower bound. A shortest route of this kind is called a \textbf{geodesic}.

The midpoint has a simpler formula: \[\boxed{g_i(1/2)=
\frac{(\sqrt{p_i}+\sqrt{r_i})^2}{2(1+B(p,r))}.}\] To derive it, normalize the vector \(\sqrt p+\sqrt r\). Its squared length is \(2+2B\), precisely the denominator. This produces the point halfway along the shorter sphere arc.

\begin{example}[Two meanings of halfway]

For \(p=(1,0)\) and \(r=(1/4,3/4)\), the coordinate midpoint is \((5/8,3/8)\). Their overlap is \(1/2\), and the spherical midpoint formula gives \((3/4,1/4)\). The latter divides the statistical distance into equal halves.

For three-category endpoints \((0.8,0.1,0.1)\) and \((0.1,0.7,0.2)\), the coordinate midpoint is \((0.45,0.4,0.15)\), while the spherical midpoint is approximately \((0.4339,0.3935,0.1726)\).

\end{example}

\begin{figure}[htbp]\centering
\resizebox{.98\linewidth}{!}{\begin{tikzpicture}\draw[ink,thick] (0,0)--(6,0)--(3,5.196)--cycle;\node[below left] at (0,0){$A$};\node[below right] at (6,0){$B$};\node[above] at (3,5.196){$C$};\draw[copper,line width=1pt,dashed] plot coordinates {(0.9,0.519615) (4.8,1.03923)};\draw[teal,line width=1pt,] plot coordinates {(0.9,0.519615) (0.979372,0.540691) (1.06171,0.561733) (1.14689,0.582706) (1.23476,0.603576) (1.32519,0.624309) (1.41801,0.64487) (1.51308,0.665225) (1.61025,0.685342) (1.70935,0.705187) (1.81022,0.724727) (1.9127,0.743931) (2.01661,0.762767) (2.12178,0.781204) (2.22805,0.799211) (2.33524,0.816759) (2.44316,0.833819) (2.55166,0.850363) (2.66053,0.866365) (2.76962,0.881796) (2.87873,0.896633) (2.98769,0.910851) (3.09632,0.924427) (3.20444,0.937337) (3.31188,0.949561) (3.41844,0.96108) (3.52397,0.971873) (3.62829,0.981923) (3.73123,0.991214) (3.83261,0.99973) (3.93227,1.00746) (4.03005,1.01438) (4.12578,1.0205) (4.21931,1.02579) (4.31049,1.03025) (4.39916,1.03387) (4.48518,1.03664) (4.56841,1.03857) (4.64871,1.03965) (4.72595,1.03987) (4.8,1.03923)};
\node[anchor=west,align=left,text=teal] at (6.6,3.8){Spherical interpolation};
\node[anchor=west,align=left,text=copper] at (6.6,3.1){Ordinary mixture};
\node[anchor=west,align=left] at (6.6,1.9){$p=(0.8,0.1,0.1)$\\$r=(0.1,0.7,0.2)$};
\end{tikzpicture}}
\caption{A spherical geodesic mapped back to the probability triangle, compared with the ordinary mixture segment between the same endpoints.}\end{figure}

In the binary case, both constructions traverse the same quarter-circle interval after mapping to the sphere, but generally at different speeds. With three categories, even the geometric routes generally differ.

For binary probabilities, equal statistical steps can be generated with particularly little machinery. Convert the starting and ending probabilities to angles \[\phi_0=\arcsin\sqrt{x_0},\qquad \phi_1=\arcsin\sqrt{x_1}.\] Choose an integer number of steps \(M\), linearly interpolate the angles, \[\phi_j=\phi_0+\frac jM(\phi_1-\phi_0),\] and convert back with \(x_j=\sin^2\phi_j\). Consecutive distances are \(2|\phi_1-\phi_0|/M\).

For the route from \(x_0=0\) to \(x_1=1\) in four steps, the angles are \[0,\quad\pi/8,\quad\pi/4,\quad3\pi/8,\quad\pi/2.\] The first probabilities are approximately \(0,0.1464,0.5,0.8536,1\). Each statistical step has length \(\pi/4\), although the probability increments are unequal.

An ordinary mixture over the same four steps gives \(0,1/4,1/2,3/4,1\). Its coordinate increments are equal, but its statistical increments are \(\pi/3,\pi/6,\pi/6,\pi/3\). Both routes traverse the same ordered binary arc, so both total lengths are \(\pi\). What differs is where the finite observations are placed along the route.

This example is useful when planning a sequence of model adjustments. Equal changes of a parameter need not mean equally spaced statistical states. A chosen coordinate system determines what ``equal step'' means.

\section{Choosing which shortest route answers the question}\label{ch08-choosing-which-shortest-route-answers-the-question}

A mixture can represent pooling sources. A spherical geodesic gives equal statistical increments. A token path respects indivisible changes. Each construction is appropriate to a different constraint.

Refining a token lattice adds available states, but adjacent denominators need not give nested lattices. Even for nested lattices, a convergence claim about shortest paths needs an argument about the allowed directions and costs. A finer picture alone is not such a proof. The exact endpoint distance remains a valid lower bound at every resolution.

In the \href{https://bobovski66.github.io/IntroDiscreteInfoGeom/demos/week08.html}{path demonstration}, compare the mixture, spherical interpolation, and permitted lattice routes for the same endpoints. Record both the endpoint distance and the sum of edge lengths. For a scalar field such as entropy, also record the signed gradient sum and verify that it depends only on the endpoints.

Chapter 9 uses geometry to choose an update rather than merely measure one. We will balance improvement in a score against the information cost of changing the current distribution.

\section{Exercises}\label{ch08-exercises}

\begin{exercises}

\item 

For potential values \(\phi(U)=0\), \(\phi(V)=2\), \(\phi(W)=5\), compute the integral on \(U,V,W,U\).

\item 

Assign one unit clockwise around a triangle. Compute both routes from \(U\) to \(W\) and explain why no potential produces the field.

\item 

Explain the root-and-path construction of a potential when every closed-path sum vanishes.

\item 

Prove the token-distance formula and construct a shortest transfer sequence from \((4,1,1)\) to \((1,2,3)\).

\item 

Verify the exact shortest length \(7\pi/6\) for the two-token example. Identify its excess over the endpoint distance.

\item 

Run Dijkstra's algorithm on the five-edge example and record the predecessor of each finalized vertex.

\item 

Explain why positive edge lengths cannot be treated as gradient increments. Use a two-edge round trip.

\item 

Verify that the ordinary mixture preserves the probability simplex.

\item 

Derive the spherical midpoint formula by normalizing \(\sqrt p+\sqrt r\).

\item 

Compute coordinate and spherical midpoints between \((1/2,1/2,0)\) and \((1/2,0,1/2)\).

\item 

Divide the binary arc from \((1,0)\) to \((1/4,3/4)\) into four equal statistical steps. Give the three interior first-coordinate values.

\item 

Choose two lattice routes with the same endpoints. Calculate a signed entropy-gradient sum and a positive statistical length for each. Explain what each total measures.

\end{exercises}

\section{Selected solutions}\label{ch08-selected-solutions}

\textbf{1.} The contributions are \(2,3,-5\), totaling zero.

\textbf{2.} The route \(U,V,W\) has sum two. The direct reverse of the clockwise edge \(W\to U\) has value \(-1\), so \(U,W\) has sum minus one. Their difference is the circulation three.

\textbf{4.} There are three surplus tokens at \(A\) and three deficits elsewhere, giving distance three. Transfer once from \(A\) to \(B\) and twice from \(A\) to \(C\).

\textbf{5.} The two possible first steps each cost \(\pi/2\). From either neighbor the direct remaining edge costs \(2\pi/3\), and no route can improve that endpoint distance. The excess is \(\pi/6\).

\textbf{10.} The coordinate midpoint is \((1/2,1/4,1/4)\). Since overlap is \(1/2\), the spherical midpoint is \((2/3,1/6,1/6)\).

\textbf{11.} The angle \(\phi=\arcsin\sqrt{x}\) decreases from \(\pi/2\) to \(\pi/6\) in four increments of \(\pi/12\). The interior first probabilities are \(\sin^2(5\pi/12)\), \(\sin^2(\pi/3)\), and \(\sin^2(\pi/4)\), or \((2+\sqrt3)/4,3/4,1/2\).
\labheading\begin{enumerate}
\def\labelenumi{\arabic{enumi}.}
\tightlist
\item
  Take a path \(A\to B\to C\) with scalar values \(\phi(A)=1\), \(\phi(B)=4\), and \(\phi(C)=2\). (a) Calculate the two gradient edge values. (b) Add them. (c) Compare with \(\phi(C)-\phi(A)\).
\item
  On a directed triangle \(A\to B\to C\to A\), assign each oriented edge value \(1\). (a) Calculate the circulation. (b) Determine whether this edge field could be the gradient of a scalar potential. (c) Justify using telescoping.
\item
  Consider count states \(a=(4,1,1)\) and \(b=(1,2,3)\) with \(N=6\). (a) Compute \(\sum_i|a_i-b_i|\). (b) Find the minimum number of token transfers. (c) Construct an explicit route attaining that minimum.
\item
  Choose two different routes between the same endpoints in the six-token triangle. (a) Calculate a scalar potential such as entropy at every vertex of both routes. (b) Sum the gradient differences. (c) Compare the unsigned cumulative changes \(\sum|\Delta H|\). (d) Explain why the latter can depend on the route.
\item
  Use Fisher--Rao edge weights for two short legal routes between the same probability states. (a) Compute and add the edge distances with a calculator. (b) Find the exact Fisher--Rao distance between endpoints. (c) Check that neither route length is smaller than the endpoint distance.
\end{enumerate}

\FloatBarrier

\end{document}
